\section{R1. Search control with coupled goals under a wall-clock budget}\label{sec:search}

\subsection{The problem and the corrected design}

The historical profile in Section~\ref{sec:profile} motivates three separate
decisions: which partial program to expand, which coupled hole to refine, and
how to allocate a finite budget among lanes and depths. In that profile one
rule opened 90\% of the tree; a depth-4 pass consumed 4.7 seconds without a
candidate while a depth-5 pass found one in 0.5 seconds. These observations
motivate experiments with resumable search, not a theorem that best-first
search or a particular cost function will improve every query.

Incoming answers N1--N9 agree on a corrected design: retain a joint native
constraint state, share persistent structure, factor only components whose
independence is justified, and separate the final candidate objective from
exploration priority and charged work. Two original requirements need repair.
``One unit per open non-leaf goal'' is not generally an admissible heuristic;
and identical candidate lists under a fixed wall-clock deadline cannot be
guaranteed across differing machine speeds while reporting every improvement
reached. The following results answer the historical R1 questions; the active
expert agenda at the end asks about the remaining implementation and research gaps.

\subsection{Conditional factorization and contextual reuse}
\label{sec:search-factorization}

There is a useful middle ground between copying all work and imposing one
global sequential order. Represent alternatives as a persistent whole-state
OR graph and factor certified independent constraints into AND components.
Genuinely incompatible assignments still require distinct logical branches.
Aesop already tracks transitive clusters of metavariable-coupled goals,
including dependencies through local contexts; independence detection itself
is not a new idea absent from Aesop~\cite{aesop}. Its copying of affected goal
nodes should not be confused with wholesale copying of every search subtree.
Canonical supplies a comparison for metavariable refinement order, and the
incoming reports identify Canonical-min's suspended constraints and
assignment-triggered resumption as another useful reference~\cite{canonicalmin}. Guarantees for
that calculus do not transfer automatically to arbitrary Lean tactics.

For each obligation, conservatively close its support through expression and
universe metavariables, assigned expressions, declaration types, local types
and values, delayed-assignment environments, postponed equations, relevant
instance arguments, provider types, and admissible local expressions. Add
residual observations as constraint nodes: $x,y:\Nat$ may have independent
types while $x+y=7$ couples their acceptable assignments. Likewise
$f(0)+g(0)=1$ couples two function bodies. An immutable shared parameter or
constant need not connect components; a writable variable in its type does.

The hypergraph can use unresolved variables as vertices and obligations as
hyperedges, or its incidence-graph dual. What matters is conservative support,
including a rule's possible writes. Treat an opaque tactic with unknown effects
as touching its entire reachable state. Overestimated dependencies lose an
optimization; underestimated ones can remove compatible completions even if
the final kernel gate rejects malformed survivors.

\begin{proposition}[Conditional component factorization]
\label{prop:search-factorization}
After fixing a shared interface assignment $\eta$, suppose the remaining
typing, grammar and behavioral constraints factor as
$\Phi_1(X_1;\eta)\land\cdots\land\Phi_k(X_k;\eta)$ with disjoint writable
sets $X_i$. Every allowed transition in component $i$ reads only $X_i$,
$\eta$ and immutable shared data, writes only $X_i$ and fresh private
variables, and preserves this separation. Then the completion set at $\eta$
is the Cartesian product of the component completion sets. Their substitutions
commute up to fresh-name renaming.
\end{proposition}
\begin{proof}
A substitution on the disjoint union of the $X_i$ is a tuple of substitutions
on its components. No right-hand side or typing dependency reads another
component's writable variables, so their capture-avoiding union is well formed.
Each conjunct is satisfied exactly when its corresponding substitution
satisfies it. Conversely, restricting a joint solution satisfies every factor.
Transition confinement ensures that this argument continues to hold after
refinement, rather than only at the initial graph snapshot.
\end{proof}

A useful operational frame condition for components $I,J$ is
\[
 W(I)\cap(R(J)\cup W(J))=\varnothing,\qquad
 W(J)\cap(R(I)\cup W(I))=\varnothing.
\]
With complete read/write footprints, immutable environment and policy, and
fresh identifiers renamed apart, each transition reads the same values in
either order and their writes unite. This is a sufficient interface contract,
not a claim that collecting visible target metavariables establishes it.
For a small unfixed separator, subsearches export relations between interface
assignments and solutions; the global result is a union of products over
compatible $\eta$. Joining those relations may still be exponential.

Update supports after assignments or context extensions. New shared helpers
or observations can merge components; assigning a shared variable can split
one. Union--find alone handles the former but not the latter. Recompute small
affected components or maintain persistent adjacency and a wakeup queue.
Rules that introduce a cross-component dependency must rejoin components or
be prohibited after a component is sealed. A global size budget, a provider
usable at most once, or a nonlinear objective is also a connecting constraint.
For additive size bounds, combine cost-indexed component solution streams;
choosing only each component's first answer can discard the only globally
admissible combination.

\paragraph{Three levels of reuse.}
Share immutable expressions and context nodes; share persistent branch prefixes
while keeping assignment maps distinct; and, only under checked interface
conditions, reuse a semantic solution. A reusable entry is a typed recipe:
abstract the minimal dependency-closed local telescope, canonicalize binders,
record the target, residual constraints, interface guard, environment,
transparency, instance policy, grammar and relevant dependency fingerprints.
On reuse freshen universe and expression variables, substitute actual locals,
validate the guard, and type-check the instantiated term. Raw \code{MVarId}s
or \code{FVarId}s cannot be transplanted into another branch. A proof of
$P(x)$ does not become a proof of $P(y)$ through similar printing.

Keep a lazy stream or guarded family of data solutions unless a checked
argument makes every remaining alternative irrelevant to the surrounding
contract. A success cache and a failure cache have different obligations:
closed contextual successes can be replayed; negative reuse requires a checked
refutation or exhaustion of a matching finite grammar and bound. Timeout,
blocked unification and a failed tactic are scheduling information, not
uninhabitability certificates.

\subsection{Admissible bounds for joint completion}
\label{sec:search-bounds}

Begin with $h(S)=0$ under fixed nonnegative costs. The original unit-per-goal
bound fails twice: exact-local closure may cost zero, and one assignment may
close several obligations. For example, assigning $?n:=0$ at cost one can
make both $n=0$ and $n+1=1$ reflexive with free propagation. Two visible
obligations then have remaining cost one. Repeating the same constraint makes
the overestimate arbitrarily large. A model charging each proof constructor
would give different costs; therefore the objective must be stated first.

Let $d(S)$ be minimum additional completion cost. If $\ell_i(S)\le d(S)$
for each justified relaxation of obligation $i$, then
\[
 h_{\max}(S)=\max_i\ell_i(S)\le d(S).
\]
A sum is valid only if every completion's cost can be partitioned into
component shares that dominate the respective bounds. Semantic independence
alone is insufficient when one provider action or shared subterm is charged
once but serves several components. Explicitly partition each concrete
action's cost without charging more than that action costs.

\begin{proposition}[Abstract-distance lower bound]
\label{prop:search-abstract-bound}
Suppose an abstraction maps concrete solved states to abstract solved states,
and every concrete transition of cost $c$ to an abstract path of cost at most
$c$. Then the shortest abstract completion distance
$d^{\#}(\alpha(S))$ is an admissible lower bound on $d(S)$.
\end{proposition}
\begin{proof}
Map a concrete completion to the concatenation of its simulating abstract
paths. It reaches an abstract solution at no greater cost. The shortest
abstract completion costs no more than this mapped path. Taking the infimum
over concrete completions gives the result.
\end{proof}

Candidate abstractions erase dependent indices, treat proof and dictionary
obligations as free, allow extra providers, retain a small separator, or use a
finite type-head production hypergraph or constructor-pattern database. A
minimum-cost cover or fractional covering relaxation can account for actions
that close several goals together. Each construction still needs the
simulation argument for \emph{every} admitted production: omitting a cast,
higher-order introduction, zero-cost local, or multi-goal unification effect
can turn a purported relaxation into an overestimate.

Derivation cost, failed search work and final-expression size are distinct.
Filling a shared hole may duplicate syntax; normalization may erase it; a
later assignment may change which inputs a completed program uses. A lower
bound on production count does not prove optimality under a different final
ranking tuple. Use a fixed additive derivation objective initially, or prove
the lower bound for the actual advertised candidate objective.

Nonnegative costs alone do not imply fair cost enumeration. Every genuinely
generative cycle needs positive cost or a finite secondary structural grade.
Administrative zero-cost closure must terminate through a well-founded
normalization argument, finite saturation or justified duplicate control;
infinitely many distinct zero-cost states evade simple cycle detection. If
an admissible but inconsistent heuristic is later used on a graph, reopen a
state when a smaller path cost is found. Uniform-cost search is the simpler
reference until a stronger bound pays for its computation.

\subsection{Learning priorities without changing the objective}
\label{sec:search-learning}

Adaptive order preserves accepted-answer correctness under the acceptance
gate, but it can change which answers appear before a deadline. Probe supplies
a just-in-time learning precedent~\cite{probe}, not a theorem of invariant
finite-budget answers or a dependent-search A* guarantee. Keep four fields
distinct: final objective $J(p)$, charged work, a proved frontier bound
$L(S)$ on unseen $J$, and dispatch priority $P_e(S)$ for epoch $e$.

Freeze $J$ for the query. Learning may reward observations survived, useful
blockers removed, estimated information gain, or low evaluation cost in
$P_e$. Correlated examples should not automatically outweigh a structurally
different observation, and training-example survival may reward overfitting.
Measure held-out behavior separately. If the advertised objective itself
changes, version it and recompute the affected bounds; silently decreasing
edge costs does not preserve the old A* argument.

Use deterministic logical epochs: collect a fixed batch of completed events,
aggregate evidence in canonical order, record the update inputs and model
digest, then freeze the next epoch's scores. Rekey affected queue entries,
rebuild the queue, or use an explicitly separate guidance queue; stale heap
keys establish neither the new order nor a reproducible old one. Canonical
ties must not depend on hash iteration, timestamps, memory addresses or
timing-derived names. Deterministic integer or rational updates are one
possible implementation, not a required statistical model.

Reserve a nonzero deterministic share of logical work for a fixed grammar
frontier. With finite branching, fair admission and terminating local work,
every finite reference derivation eventually receives service. A beam or
top-$k$ list in a fast lane may be useful, but omitted alternatives must remain
scheduled in the reference lane for a coverage claim. Introducing a helper
or carrier outside the grammar is a recorded language extension, not merely
reordering. Exhaustion of the old grammar is not exhaustion of the enlarged one.

The appropriate default receipt is ``best certified candidates found under
objective version $v$ in this logical prefix.'' A stronger minimum-cost claim
requires an incumbent $p$ and coverage of \emph{all} unexplored alternatives
under the same objective, with
\[
 J(p)\le\inf_{S\text{ on the complete frontier}}L(S).
\]
A grace period after the first answer is not such a certificate. Stable
sorting gives stable order of a given verified set, not stability of the set
that a time-limited run discovers. Test user-facing stability under irrelevant
renaming, harmless provider permutations, repeated runs and additional
examples rather than treating an algorithm name as evidence of predictability.

\subsection{Deterministic prefixes and real deadlines}
\label{sec:search-deadlines}

A reproducible logical schedule and a hard deadline are compatible, but the
deadline may truncate that schedule at different positions. The anytime
literature supplies a quality-versus-deliberation perspective~\cite{anytime};
it does not establish machine-independent equality of timed answer sets.

\begin{proposition}[Deadline obstruction]\label{prop:search-deadline}
Suppose a deterministic search first certifies a new answer after positive
work $w$. If one execution reaches $w$ before deadline $D$ and another does
not, an interface publishing every certified improvement reached by $D$
cannot return the same answer set in both executions.
\end{proposition}
\begin{proof}
The answer is available only in the first completed prefix. Sorting cannot
insert an undiscovered certified answer into the other. Suppressing the new
answer or continuing beyond the deadline changes one of the stated requirements.
\end{proof}

Even a restricted speed variation breaks equality if a cutoff crosses an
answer event. A known throughput bound can justify fixed work for a specified
machine class; startup calibration alone cannot account for cache warmth,
instance search, reduction costs, load or allocation. Nor does one expensive
node represent a stable amount of work.

Provide an interactive execution contract that respects cancellation and
reports its completed prefix, a replay contract that consumes that prefix
regardless of elapsed time, and fixed-work runs for experiments. Pin query,
environment, grammar, policy, provider order, model, seed and tie-breaks.
Checkpoints contain the frontier, semantic branch states, expansion cursors,
model state, monotone work ledger, unresolved tasks and exhausted cost bands,
not only a candidate count. Charge rule, unification, normalization, instance,
observation and proof work at specified points even after rollback. Resuming
to the same logical checkpoint should reproduce the stream; three timed
two-second runs need not equal one six-second run because startup and replay
also consume time.

Parallel workers may speculate, but stable task identities and fixed logical
allotments must feed a deterministic commit order. Publish completed epochs;
label any opportunistic extras separately. A race to the first answer followed
by cancellation is not repaired by sorting its survivors. Cooperative node
checks and heartbeats are not universal real-time interrupts: bound expensive
services and use a worker cancellation boundary where necessary. Reserve
resources for final checking; interruption cannot promote an unchecked term.

Iterative deepening remains a useful baseline: its familiar large-branching
overhead argument concerns completed brute-force passes~\cite{korf85}, not
arbitrarily cut heterogeneous lanes. Likewise the Luby restart theorem assumes
a Las Vegas runtime model~\cite{luby93}; correlated deterministic lanes with
shared learned state do not automatically meet it. Prefer resumable work to
restarting the identical prefix. Compare restarts only when a new policy or
random trial gives them a substantive purpose.

\subsection{Implementation sequence and evidence still needed}
\label{sec:search-experiments}

First make the \emph{whole semantic state} and a resumable expansion cursor
explicit while reproducing the continuation baseline with a LIFO frontier.
\code{Meta.SavedState} alone is not a specification of every branch-dependent
cache or elaborator effect. Audit external references and residual-cache
identities; semantic caches need valid branch keys, while global work charges
remain outside rollback. Snapshot creation cost does not measure the memory
retained by a broad persistent frontier. External process snapshotting and
in-process persistent contexts have different cost regimes. N5's supplied
snapshotting study is a relevant tactic-search precedent~\cite{snapshot},
not evidence for the cost of this proposed in-process frontier.

Next use the dependency graph for goal pressure and blocker wakeups, then
memoize closed contextual successes, then introduce certified independent
commits and separator tables if repeated work justifies them. Fail-first
selection can estimate the applicable rule domain from retrieval, with ties
favoring assignments that constrain other goals, and aging prevents starvation.
Agsy and Canonical are useful comparisons~\cite{agsy,canonical}; a dependent
goal's domain is an estimate, not an exact finite set supplied by those systems.
The historical initial weights (exact 0, constructor/local application 1,
provider 2, split 3, recursion 4, accumulator choice 5) remain hypotheses to
measure under the zero-cost safeguards, not a proved optimal cost model.

A bottom-up companion remains a distinct option: enumerate small typed terms
over observation environments and offer components to top-down holes, following
the broad ideas of Escher and Burst~\cite{escher,burst}. Finite observational
signatures may organize buckets, but retaining only one representative needs
a contextual congruence argument for the admitted grammar or a way to recover
discarded alternatives. Observational agreement is not unrestricted program
equality. Grammar learning and divide-and-conquer enumeration remain useful
comparisons~\cite{deepcoder,eusolver,transit}, not replacements for these checks.

Measure solved count at 1, 10 and 60 seconds, first certified answer, best fixed
objective at equal logical work, proof debt, frontier memory, repeated work,
component sizes, cache replay cost and nodes off the successful path. Compare
the same grammar with and without goal ordering, factorization, heuristics
and learning. Replace the impossible demand of zero answer-set variance under
$\pm30\%$ speed changes with equal outcomes at equal logical checkpoints;
report time-dependent prefix differences and actual latency separately.
The incoming Python countermodels illustrate failed bounds and deadline
claims. They are not native Lean validation or performance measurements.

\subsection{New questions for experts}\label{sec:search-new-questions}
The historical questions have the conditional answers above. The next questions
concern the missing constructive analyses and matched experiments.
\begin{description}
\item[R1.N1.] What conservative read/write analysis covers Lean's delayed
assignments, local instances and elaborator effects cheaply enough to justify
component factorization, and which workloads repay separator joins and retained
frontier memory?
\item[R1.N2.] Which finite relaxation admits a rule-by-rule cost simulation for
the actual final candidate objective, including sharing and normalization, while
providing enough pruning to repay its computation?
\item[R1.N3.] Which deterministic learning signals improve held-out synthesis
at equal logical work without destabilizing fixed-objective ranking under
irrelevant renaming, provider-order changes and correlated examples?
\item[R1.N4.] What resumable checkpoint and interruption protocol can reproduce
native search, proof services and parallel proposal admission at equal logical
prefixes with acceptable latency and memory overhead?
\end{description}
